%
% Chapitre "Structure d'un rapport technique"
%

\chapter{Étude préliminaire}
\label{s:etude_preliminaire}
\section{Plan de développement}
Afin d'asurer une évaluation systématique de chaque concept, un plan d'étude a été élaboré.
Le tableau \ref{tab:plan_etude} présente en détail la procédure d'évaluation de chacun des critères, appliqué aux trois concepts exposés dans le tableau \ref{tab:concepts_nordfaune}.


\newcolumntype{L}[1]{>{\RaggedRight\arraybackslash}p{#1}}
\renewcommand{\arraystretch}{1.4}

\begin{longtable}{|L{4cm}|L{6cm}|L{6cm}|}
\caption{Plan d'étude pour les différents critères}
\label{tab:plan_etude} \\

\hline
\textbf{Critères} & \textbf{Procédures} & \textbf{Hypothèses} \\
\hline
\endfirsthead

\multicolumn{3}{c}{\textit{Table 6.1 -- (suite)}} \\
\hline
\textbf{Critères} & \textbf{Procédures} & \textbf{Hypothèses} \\
\hline
\endhead

% ------------- LIGNES -------------

\ref{robutesse_autonomie} Marge de l'état d'autonomie
    & Évaluer l'autonomie la plus basse du système selon les données du fournisseur.
    & Toutes les composantes fonctionnent dans leurs plages nominales.
\\
\hline

\ref{consommation_electrique} Consommation énergétique
    & Évaluer la somme des consommation électrique selon les données du fournisseur.
    & Le système fonctionne de manière continue sous conditions climatiques normales.
\\
\hline

\ref{resilience_systeme} Proportion des fonctions protégées
    & Évaluer le nombre de fonctions protégées par la redondance ou une mesure de protection.
    & Chaque fonction critique peut être protégée indépendamment.
\\
\hline

\ref{conditions_climatiques} Plage opérationnelle de température
    & Évaluer la température minimale d'opération selon les données du fournisseur.
    & L'isolation thermique du système est uniforme et stable.
\\
\hline

\ref{installation_deploiement} Encombrement et complexité d'installation 
    & Évaluer le volume du sytème et sa manipulation selon les données du fournisseur.
    & La forme du système est manipulable par deux personnes.
\\
\hline

\ref{resistance_neige} Résistance mécanique à la neige 
    & Évaluer la charge maximale supportée par le boîtier selon les données du fournisseur.
    & La neige est répartie uniformément sur le boîtier.
\\
\hline

\ref{detection_mammiferes} Couverture de la zone d'observation
    & Évaluer par projection géométrique la proportion de l'aire couverte.
    & Les capteurs ne sont pas obstrués et sont placés à l'endroit optimal.
\\
\hline

\ref{conservation_donnees} Capacité de stockage
    & Évaluer à partir de l'encodage, la fréquence, la durée et la résolution de la capture.
    & La compression et le format sont constants et uniformes.
\\
\hline

\ref{performance_communication} Performance de communication distante
    & Évaluer la portée maximale selon les données du fournisseur.
    & La ligne de communication n'a aucun obstacles majeurs.
\\
\hline

\ref{exploitation_scientifique} Qualité scientifique des données
    & Évaluer la fréquence d'image que l'appareil de capture offre selon le fabricant.
    & Les capteurs vidéo fonctionne à cadence régulière.
\\
\hline

\ref{impact_lumineux} Impact lumineux sur la faune
    & Évaluer avec la longueur d'onde de la lumière visible produite.
    & Les lemmings ne sont sensibles qu'à la lumière dans la plage définie.
\\
\hline

\ref{impact_environnemental} Impact environnemental
    & Évaluer le nombre de batteries consommées annuellement selon le fournisseur.
    & Chaque batterie a une autonomie prévisible.
\\
\hline

\ref{couts_transport} Coûts d'installation et de transport
    & Évaluer en fonction de la somme des frais d'installation des éléments cruciaux.
    & Les coûts de transport et d'installation sont proportionnels au temps exigé.
\\
\hline

\ref{couts_entretien} Coûts d'opération et d'entretien
    & Évaluer selon le coûts du renouvellement des mesures de protection, les frais qu'apporte l'opération et d'autres coûts obligatoires.
    & Les frais récurrents sont stables.
\\
\hline

\ref{couts_fabrication} Coûts de fabrication
    & Évaluer avec la somme des coûts des composantes majeures à l'achat.
    & Les prix des composants sont constants et correspondent aux prix affichés.
\\
\hline

\end{longtable}


\renewcommand{\arraystretch}{1.5}

\begin{table}[htp]
\centering
\caption{Les trois concepts pour le projet NordFaune}
\label{tab:concepts_nordfaune}

\begin{tabular}{|L{6cm}|L{3cm}|L{3cm}|L{3cm}|}
\hline
\textbf{Fonctions} & \textbf{Concept 1} & \textbf{Concept 2} & \textbf{Concept 3} \\
\hline

% ------------- LIGNES -------------

\ref{CapturerUneVidéo} Capturer une vidéo
    & \hyperref[sec:arduc]{ArduCam IMX219}
    & \hyperref[sec:gopro]{GoPro Hero}
    & \hyperref[sec:esp32]{ESP32-CAM}
\\
\hline

\ref{DétecterLeMouvement} Détecter le mouvement
    & \hyperref[sec:plaque]{Plaque piézoélectriques}
    & \hyperref[sec:capte]{Capteur infrarouge}
    & \hyperref[sec:capte]{Capteur infrarouge}
\\
\hline

\ref{ContrôlerLAccès} Contrôler l'accès
    & \hyperref[sec:tunne]{Tunnel Cloudflare}
    & \hyperref[sec:micro]{Microsoft Authentificator}
    & \hyperref[sec:tunne]{Tunnel Cloudflare}
\\
\hline

\ref{EnregistrerDesDonnées} Enregistrer des données
    & \hyperref[sec:cartesd]{MicroSD(HC)}
    & \hyperref[sec:ssd]{SSD SATA III}
    & \hyperref[sec:azsta]{Azure et Starlink}
\\
\hline

\ref{DétecterUnÉnévement} Détecter un événement et l'état du système
    & \hyperref[sec:raspb]{Raspberry Pi}
    & \hyperref[sec:nano]{Nvidia Jetson Nano}
    & \hyperref[sec:raspb]{Raspberry Pi}
\\
\hline

\ref{RésisterAuxConditions} Résister aux conditions climatiques
    & \hyperref[sec:bois]{Bois de cèdre}
    & \hyperref[sec:alumi]{Aluminium}
    & \hyperref[sec:alumi]{Aluminium}
\\
\hline

\ref{GérerLÉnergie} Gérer l'énergie
    & \hyperref[sec:bloc]{Bloc de batteries}
    & \hyperref[sec:hybri]{Solution hybride}
    & \hyperref[sec:bloc]{Bloc de batteries}
\\
\hline

\ref{CommuniquerAvecExt} Communiquer avec l'extérieur
    & \hyperref[sec:iridi]{Iridium 9603 SBD}
    & \hyperref[sec:starl]{Starlink Performance}
    & \hyperref[sec:iridi]{Iridium 9603 SBD}
\\

\hline
\end{tabular}
\end{table}


\section{Élaboration et évaluation des concepts de solution}

% ------------- Concept #1 -------------

\subsection{Concept \#1 : Équilibre global}

Pour ce concept, les composantes ont été choisies de façon à garder un bon équilibre dans l'ensemble des fonctions du système, sans trop dépenser.

\subsubsection{Fonctionnement autonome du système}

\paragraph{Marge de l'état d'autonomie :}
Étant donné que ce concept utilise un contrôleur Raspberry Pi, une caméra ArduCam IMX219 et un module Iridium 9603 SBD, la consommation du système demeure modérée.
Le système se repose complètement sur les blocs de batteries non rechargeables, qui vont s’épuiser avant la durée de vie des composants électroniques.

\bigskip
La contrainte du client exige une autonomie minimale de douze mois.
Il est entièrement possible d’atteindre cette attente avec une quantité de batteries proportionnelle à la consommation.
Le facteur limitant est donc la capacité totale des batteries embarquées.
Alors, avec un dimensionnement optimisé (Voir \hyperref[par:encombrement]{Encombrement}), l’autonomie maximale du concept est estimée à environ 12,5 mois.
D’après le barème de la section \ref{robutesse_autonomie}, cette solution obtient une note de 0,29 pour ce concept, soit 2,9\% avec pondération.


\paragraph{Consommation énergétique :}
Texte

\paragraph{Proportion des fonctions protégées :}
Texte


\subsubsection{Résistance et entretien du système}

\paragraph{Plage opérationnelle de température :}
Texte

\paragraph{Encombrement et complexité d'installation :}
\label{par:encombrement}
Texte

\paragraph{Résistance mécanique à la neige :}
Texte


\subsubsection{Observation et gestion de données}

\paragraph{Couverture de la zone d'observation :}
Texte

\paragraph{Capacité de stockage :}
Texte

\paragraph{Performance de communication distante :}
Texte

\paragraph{Qualité scientifique des données :}
Texte


\subsubsection{Réduire impacts et coûts}

\paragraph{Impact lumineux sur la faune :}
Texte

\paragraph{Impact environnemental :}
Texte

\paragraph{Coûts d'installation et de transport :}
Texte

\paragraph{Coûts d'opération et d'entretien :}
Texte

\paragraph{Coûts de fabrication :}
Texte

% ------------- Concept #2 -------------

\subsection{Concept \#2 : Performance maximale}

Pour ce concept, les options les plus performantes ont été privilégiées pour garantir la meilleure performance, peu importe le coût.

\subsubsection{Fonctionnement autonome du système}

\paragraph{Marge de l'état d'autonomie :}
L’obtention d’une performance maximale apporte des composants très énergivores.
Une GoPro, un Nvidia Jetson Nano et un système Starlink ont une demande énergétique trop grande pour dépendre sur un système énergétique standard.
Alors, le concept exige la solution hybride combinant des batteries au Lithium-Thionyl Chloride et des condensateurs hybrides au lithium.
\bigskip

Le facteur limitant devient la durée maximale de vie planifiée lors de la conception des batteries $\mathrm{Li{-}SOCl_2}$.
Dans un circuit soumis à une très forte utilisation, elles sont capables d’atteindre une durée de 12 mois.
Nous estimons que, pour notre système qui utilise moins d’énergie que la valeur nominale, l’autonomie peut être assurée jusqu’ à 14 mois indépendamment des conditions externes.
Selon l’équation \ref{robutesse_autonomie}, le concept obtient donc une note de 0,79, ce qui correspond à un pourcentage de 7,9\%.


\paragraph{Consommation énergétique :}
Texte

\paragraph{Proportion des fonctions protégées :}
Texte


\subsubsection{Résistance et entretien du système}

\paragraph{Plage opérationnelle de température :}
Texte

\paragraph{Encombrement et complexité d'installation :}
Texte

\paragraph{Résistance mécanique à la neige :}
Texte


\subsubsection{Observation et gestion de données}

\paragraph{Couverture de la zone d'observation :}
Texte

\paragraph{Capacité de stockage :}
Texte

\paragraph{Performance de communication distante :}
Texte

\paragraph{Qualité scientifique des données :}
Texte


\subsubsection{Réduire impacts et coûts}

\paragraph{Impact lumineux sur la faune :}
Texte

\paragraph{Impact environnemental :}
Texte

\paragraph{Coûts d'installation et de transport :}
Texte

\paragraph{Coûts d'opération et d'entretien :}
Texte

\paragraph{Coûts de fabrication :}
Texte

% ------------- Concept #3 -------------

\subsection{Concept \#3 : Durabilité à long terme}

Pour ce concept, les solutions les plus durables ont été sélectionnées afin d'assurer un bon fonctionnement à long terme.

\subsubsection{Fonctionnement autonome du système}

\paragraph{Marge de l'état d'autonomie :}
Étant donné que le concept est composé d’une ESP32-CAM, d’un Raspberry Pi et d’un module Iridium 9603 SBD, la demande énergétique du système est fortement réduite.
L’alimentation est assurée par un bloc de batteries optimisé.
\bigskip

Dans ce cas, le facteur limitant est plutôt la durée de vie maximale des composants électroniques, notamment le module ESP32-CAM et le module Iridium, qui ont une fiabilité réduite lors d’une utilisation sur une longue période dans des conditions extrêmes.
L’autonomie est estimée comme étant assurée pendant 18 mois.
Cette durée lui donne une note de 1 selon l’équation \ref{robutesse_autonomie}, ce qui correspond à un pourcentage de 10\%.


\paragraph{Consommation énergétique :}
Texte

\paragraph{Proportion des fonctions protégées :}
Texte


\subsubsection{Résistance et entretien du système}

\paragraph{Plage opérationnelle de température :}
Texte

\paragraph{Encombrement et complexité d'installation :}
Texte

\paragraph{Résistance mécanique à la neige :}
Texte


\subsubsection{Observation et gestion de données}

\paragraph{Couverture de la zone d'observation :}
Texte

\paragraph{Capacité de stockage :}
Texte

\paragraph{Performance de communication distante :}
Texte

\paragraph{Qualité scientifique des données :}
Texte


\subsubsection{Réduire impacts et coûts}

\paragraph{Impact lumineux sur la faune :}
Texte

\paragraph{Impact environnemental :}
Texte

\paragraph{Coûts d'installation et de transport :}
Texte

\paragraph{Coûts d'opération et d'entretien :}
Texte

\paragraph{Coûts de fabrication :}
Texte



\section{Synthèse des résultats}
Le tableau \ref{tab:synth} présente les valeurs obtenues pour chaque critère en fonction des différents concepts.
Ces valeurs permettent ensuite de calculer la note obtenue à l'aide des barèmes du \hyperref[s:cahier_des_charges]{cahier des charges}.

\begin{table}[h!]
\centering
\caption{Synthèse des résultats}
\label{tab:synth}
\renewcommand{\arraystretch}{1.2}
\begin{tabularx}{\textwidth}{|X|c|c|c|}
\hline
\textbf{Critères d'évaluation} & \textbf{Concept 1} & \textbf{Concept 2} & \textbf{Concept 3} \\
\hline

\textbf{\ref{Autonomie} Fonctionnement autonome du système} & & &  \\
\hline
\ref{robutesse_autonomie} Marge de l'état d'autonomie [mois] & 12,5 & 14 & 18 \\
\ref{consommation_electrique} Consommation énergétique [W] & & & \\
\ref{resilience_systeme} Proportion des fonctions protégées & & & \\
\hline

\textbf{\ref{Robustesse} Résistance et entretien du système} & & &  \\
\hline
\ref{conditions_climatiques} Plage opérationnelle de température [°C] & & & \\
\ref{installation_deploiement} Encombrement et complexité d’installation [$m^3$] & & & \\
\ref{resistance_neige} Résistance mécanique à la neige [$kg/m^2$] & & & \\
\hline

\textbf{\ref{donnees} Observation et gestion de données} & & & \\
\hline
\ref{detection_mammiferes} Couverture de la zone d’observation [\%] & & & \\
\ref{conservation_donnees} Capacité de stockage [Go] & & & \\
\ref{performance_communication} Performance de communication distante [km] & & & \\
\ref{exploitation_scientifique} Qualité scientifique des données [fps] & & & \\
\hline

\textbf{\ref{impacts_couts} Réduire impacts et coûts} & & &  \\
\hline
\ref{impact_lumineux} Impact lumineux sur la faune [nm] & & & \\
\ref{impact_environnemental} Impact environnemental & & & \\
\ref{couts_transport} Coûts d'installation et de transport [\$] & & & \\
\ref{couts_entretien} Coûts d'opération et d’entretien [\$] & & & \\
\ref{couts_fabrication} Coûts de fabrication [\$] & & & \\
\hline

\end{tabularx}
\end{table}
